% =============================================================================
% Resumen teorico - Practico 01 (Maquinas de Turing)
% Computabilidad y Complejidad (Teoria de la Computacion II)
% FCEFQyN - Universidad Nacional de Rio Cuarto
% Primer Cuatrimestre de 2026
% -----------------------------------------------------------------------------
% Fuentes principales:
%   - chapter-03.pdf   (Sipser, Cap. 3: The Church-Turing Thesis)
%   - MaqTuring.pdf    (laminas de la catedra: Maquinas de Turing)
%   - Diag.pdf         (laminas de la catedra: Indecidibilidad / Cantor)
%
% Compilacion sugerida: pdflatex teoria-practico-01.tex
% =============================================================================

\documentclass[11pt,a4paper]{article}

% --- Codificacion y idioma --------------------------------------------------
\usepackage[utf8]{inputenc}
\usepackage[T1]{fontenc}
\usepackage[spanish,es-tabla]{babel}
\usepackage{lmodern}
\usepackage{microtype}

% --- Matematica -------------------------------------------------------------
\usepackage{amsmath,amssymb,amsthm,mathtools}
\usepackage{stmaryrd}

% --- Layout y tipografia ----------------------------------------------------
\usepackage[margin=2.7cm]{geometry}
\usepackage{enumitem}
\setlist[itemize]{topsep=2pt,itemsep=2pt,parsep=0pt}
\setlist[enumerate]{topsep=2pt,itemsep=2pt,parsep=0pt}
\usepackage{parskip}
\usepackage{xcolor}
\usepackage{booktabs}
\usepackage{array}

% --- Hyperref ---------------------------------------------------------------
\usepackage[hidelinks]{hyperref}
\hypersetup{
  pdftitle={Resumen teorico - Practico 01},
  pdfauthor={Computabilidad y Complejidad - UNRC 2026},
  pdfsubject={Maquinas de Turing - Resumen teorico},
  pdfkeywords={Maquina de Turing, Sipser, Church-Turing, decidibilidad, reconocibilidad}
}

% --- Ambientes de teoremas --------------------------------------------------
\theoremstyle{plain}
\newtheorem{teorema}{Teorema}[section]
\newtheorem{proposicion}[teorema]{Proposicion}
\newtheorem{lema}[teorema]{Lema}
\newtheorem{corolario}[teorema]{Corolario}

\theoremstyle{definition}
\newtheorem{definicion}[teorema]{Definicion}
\newtheorem{ejemplo}[teorema]{Ejemplo}

\theoremstyle{remark}
\newtheorem*{observacion}{Observacion}
\newtheorem*{intuicion}{Intuicion}
\newtheorem*{idea}{Idea}

% --- Comandos de notacion ---------------------------------------------------
\newcommand{\MT}{\mathrm{MT}}
\newcommand{\TM}{\mathrm{TM}}
\newcommand{\enc}[1]{\langle #1\rangle}
\newcommand{\lang}{\mathcal{L}}
\newcommand{\Sstar}{\Sigma^{*}}
\newcommand{\qaccept}{q_{\text{accept}}}
\newcommand{\qreject}{q_{\text{reject}}}
\newcommand{\blank}{\sqcup}
\newcommand{\TMstd}{\TM_{\text{std}}}
\newcommand{\TMtwo}{\TM_{2}}
\newcommand{\TMone}{\TM_{1}}
\newcommand{\TMk}{\TM_{k}}
\newcommand{\NTM}{\mathrm{NTM}}
\newcommand{\rec}{\mathrm{REC}}
\newcommand{\dec}{\mathrm{DEC}}
\newcommand{\ATM}{A_{\TM}}
\newcommand{\HALT}{\mathit{HALT}_{\TM}}
\newcommand{\Pow}{\mathcal{P}}

% --- Cabecera ---------------------------------------------------------------
\title{\textbf{Resumen Te\'orico -- Pr\'actico 01}\\[0.3em]
       \large M\'aquinas de Turing\\[0.4em]
       \normalsize Computabilidad y Complejidad (Teor\'ia de la Computaci\'on II)\\
       \normalsize FCEFQyN -- Universidad Nacional de R\'io Cuarto\\
       \normalsize Primer Cuatrimestre de 2026}
\author{}
\date{}

\begin{document}
\maketitle
\thispagestyle{empty}

\hrule
\vspace{0.4em}
\begin{center}
\small
S\'intesis de la teor\'ia necesaria para resolver el Pr\'actico 01.
Apoyada en el Cap\'itulo 3 del libro de M.~Sipser,
\emph{Introduction to the Theory of Computation}, y en las l\'aminas
de la c\'atedra (\texttt{MaqTuring.pdf}, \texttt{Diag.pdf}).
\end{center}
\vspace{0.4em}
\hrule

\tableofcontents
\bigskip
\hrule
\bigskip

% =============================================================================
\section{Introducci\'on y motivaci\'on hist\'orica}
\label{sec:intro}

Las m\'aquinas de Turing (MT) son la formalizaci\'on cl\'asica de la noci\'on
intuitiva de \emph{algoritmo}. Aparecieron en el art\'iculo de Alan~Turing
de 1936, \emph{On Computable Numbers, with an Application to the
Entscheidungsproblem}, motivado por el d\'ecimo problema de la lista que
Hilbert hab\'ia planteado en 1900.

\paragraph{El d\'ecimo problema de Hilbert.} El problema preguntaba si
\emph{existe un procedimiento finito} para determinar si un polinomio con
coeficientes enteros tiene una ra\'iz integral. Por ejemplo, si se da el
polinomio $6x^3yz^2 + 3xz^2$, decidir si existen $x,y,z\in\mathbb Z$ tales
que el polinomio se anula. El problema en s\'i mismo es matem\'atico, pero
exhibe una dificultad meta-l\'ogica: para responder ``no existe tal
procedimiento'' necesitamos primero \emph{definir formalmente} qu\'e es un
``procedimiento''. Es decir, hace falta una definici\'on precisa de la
noci\'on intuitiva de algoritmo.

\paragraph{La respuesta de Turing.} Turing propuso definir un algoritmo
como una m\'aquina abstracta sumamente simple --- la que hoy llamamos
\emph{m\'aquina de Turing} --- y demostr\'o que ya con esa formalizaci\'on
hay problemas que ning\'un procedimiento finito puede resolver. El d\'ecimo
problema de Hilbert resulta ser uno de ellos (Matij\'asevich, 1970).

\begin{intuicion}
Antes de Turing, el concepto de ``algoritmo'' era informal: la criba de
Erat\'ostenes o el algoritmo de Euclides eran procedimientos descriptos en
lenguaje natural. Esta informalidad bastaba para \emph{exhibir} algoritmos,
pero era in\'util para \emph{demostrar que un problema no admite ninguno}.
Para probar la imposibilidad de un algoritmo hace falta una definici\'on
matem\'atica de lo que es un algoritmo. La MT cumple ese rol.
\end{intuicion}

\paragraph{Por qu\'e este Pr\'actico.} El Pr\'actico 01 introduce todas las
piezas necesarias para hablar de algoritmos formalmente: la definici\'on de
MT, sus configuraciones y c\'omputos, sus variantes principales (cinta
doblemente infinita, multicinta, no determinista, enumeradores), y los
conceptos de \emph{decidibilidad} y \emph{reconocibilidad}. Este resumen
recorre cada uno de esos temas con suficiente profundidad para resolver y
entender los ejercicios.

\bigskip\hrule\bigskip

% =============================================================================
\section{Definici\'on formal de m\'aquina de Turing}
\label{sec:formal}

\subsection{Componentes de una MT}

Intuitivamente, una m\'aquina de Turing consiste en:

\begin{itemize}
  \item una \emph{cinta} infinita dividida en celdas, cada una conteniendo
        un s\'imbolo;
  \item un \emph{cabezal} que apunta a una celda y puede leer y escribir;
  \item un \emph{control finito} que, en funci\'on del estado actual y del
        s\'imbolo le\'ido, decide qu\'e escribir, hacia d\'onde mover el
        cabezal y a qu\'e estado pasar.
\end{itemize}

A diferencia de un aut\'omata finito, la MT \emph{escribe} en su cinta y
\emph{recorre} la cinta en ambos sentidos; a diferencia de un aut\'omata
con pila, su memoria es \emph{ilimitada y de acceso libre}.

\begin{definicion}[7-tupla, Sipser 3.3]
\label{def:mt}
Una m\'aquina de Turing es una 7-tupla
\[
  M = (Q, \Sigma, \Gamma, \delta, q_0, \qaccept, \qreject),
\]
donde $Q,\Sigma,\Gamma$ son conjuntos finitos y:
\begin{enumerate}
  \item $Q$ es el conjunto de \emph{estados};
  \item $\Sigma$ es el \emph{alfabeto de entrada}, con $\blank\notin\Sigma$;
  \item $\Gamma$ es el \emph{alfabeto de cinta}, con $\Sigma\cup\{\blank\}\subseteq\Gamma$;
  \item $\delta : Q\times\Gamma \to Q\times\Gamma\times\{L,R\}$ es la
        \emph{funci\'on de transici\'on};
  \item $q_0\in Q$ es el \emph{estado inicial};
  \item $\qaccept\in Q$ es el \emph{estado de aceptaci\'on};
  \item $\qreject\in Q$ es el \emph{estado de rechazo}, con $\qaccept\neq\qreject$.
\end{enumerate}
\end{definicion}

El s\'imbolo $\blank$ se llama \emph{blanco} y representa una celda vac\'ia
de la cinta. Los s\'imbolos en $\Gamma\setminus\Sigma$ son auxiliares y
suelen usarse para marcar celdas durante la computaci\'on (cruces, marcas,
delimitadores).

\begin{intuicion}
Pensar a $\delta(q,a)=(p,b,D)$ como la regla \emph{``si estoy en $q$ y leo
$a$, entonces escribo $b$, me muevo en direcci\'on $D$ y paso al estado
$p$''}. La direcci\'on $D\in\{L,R\}$ es izquierda o derecha.
\end{intuicion}

\subsection{Configuraciones}

Para describir el estado din\'amico de la m\'aquina (qu\'e tiene en la cinta,
en qu\'e estado est\'a, donde apunta el cabezal) se usa el concepto de
\emph{configuraci\'on}.

\begin{definicion}[Configuraci\'on]
\label{def:configuracion}
Una \emph{configuraci\'on} de $M$ es una cadena de la forma
\[
  u\,q\,v,\qquad u,v\in\Gamma^*,\ q\in Q,
\]
con la siguiente interpretaci\'on:
\begin{itemize}
  \item el contenido de la cinta es $uv$ seguido de blancos;
  \item el cabezal est\'a sobre el primer s\'imbolo de $v$;
  \item el estado actual es $q$.
\end{itemize}
\end{definicion}

\begin{ejemplo}
La configuraci\'on $1011 q_7 01111$ significa: el contenido de la cinta es
$1011\,01111$, el cabezal est\'a sobre el primer $0$ de $01111$, y el estado
actual es $q_7$. Esquem\'aticamente:
\[
  \boxed{1}\,\boxed{0}\,\boxed{1}\,\boxed{1}\,\boxed{\underset{\uparrow q_7}{0}}\,\boxed{1}\,\boxed{1}\,\boxed{1}\,\boxed{1}\,\blank\,\blank\,\cdots
\]
\end{ejemplo}

\subsection{La relaci\'on ``cede a''}

Las configuraciones se relacionan por la regla de transici\'on:

\begin{definicion}[Yields, $\to_M$]
\label{def:yields}
Sean $C$ y $C'$ dos configuraciones de $M$. Decimos que $C$ \emph{cede a}
$C'$ en $M$, notado $C\to_M C'$, si $C'$ se obtiene de $C$ aplicando una
transici\'on de $\delta$. Mas formalmente, sea $a,b,c\in\Gamma$,
$u,v\in\Gamma^*$ y $q_i,q_j\in Q$.
\begin{itemize}
  \item Si $\delta(q_i,b)=(q_j,c,L)$, entonces
        \(u a q_i b v \;\to_M\; u q_j a c v.\)
  \item Si $\delta(q_i,b)=(q_j,c,R)$, entonces
        \(u q_i b v \;\to_M\; u c q_j v.\)
  \item En el extremo izquierdo, $q_i b v \to_M q_j c v$ si la transici\'on
        es a la izquierda (el cabezal no se mueve, queda en su lugar).
  \item En el extremo derecho, $u q_i$ es equivalente a $u q_i \blank$ y se
        trata como en los casos anteriores.
\end{itemize}
\end{definicion}

\subsection{Configuraci\'on inicial, final, c\'omputo}

\begin{definicion}[Configuraci\'on inicial]
La configuraci\'on inicial de $M$ con entrada $w=w_1w_2\cdots w_n$ es
$q_0w$, es decir: la cinta contiene $w$ en sus primeras $n$ celdas, el
cabezal apunta al primer s\'imbolo, el estado es $q_0$.
\end{definicion}

\begin{definicion}[Configuraci\'on de aceptaci\'on / rechazo]
Una configuraci\'on $u\,q\,v$ es de \emph{aceptaci\'on} si $q=\qaccept$ y de
\emph{rechazo} si $q=\qreject$. Ambas son \emph{configuraciones de
detenci\'on} ($M$ no ejecuta m\'as transiciones).
\end{definicion}

\begin{definicion}[$M$ acepta $w$]
\label{def:acepta}
$M$ \emph{acepta} la cadena $w$ si existe una sucesi\'on finita de
configuraciones $C_0,C_1,\dots,C_k$ tal que:
\begin{enumerate}
  \item $C_0 = q_0 w$ (configuraci\'on inicial);
  \item $C_i \to_M C_{i+1}$ para todo $0\le i < k$;
  \item $C_k$ es una configuraci\'on de aceptaci\'on.
\end{enumerate}
La sucesi\'on $C_0,\dots,C_k$ se llama \emph{c\'omputo} (o \emph{run})
\emph{aceptante} de $M$ sobre $w$.
\end{definicion}

\subsection{Lenguaje reconocido por una MT}

\begin{definicion}[Lenguaje de $M$]
El \emph{lenguaje reconocido} por $M$ es
\[
  \lang(M) \;=\; \{\,w\in\Sstar : M \text{ acepta } w\,\}.
\]
\end{definicion}

\bigskip\hrule\bigskip

% =============================================================================
\section{Reconocibilidad y decidibilidad}
\label{sec:rec-dec}

Las MT no siempre se detienen: pueden caer en un bucle. Distinguimos dos
clases fundamentales de lenguajes seg\'un el comportamiento de la MT que los
caracteriza.

\subsection{Tres resultados posibles}

Cuando $M$ se ejecuta sobre una entrada $w$, hay exactamente \emph{tres}
desenlaces posibles:

\begin{itemize}
  \item $M$ \emph{acepta} $w$ (alcanza $\qaccept$).
  \item $M$ \emph{rechaza} $w$ (alcanza $\qreject$).
  \item $M$ \emph{cicla} (no alcanza ning\'un estado de detenci\'on, su
        c\'omputo es infinito). Tambi\'en se dice que $M$ \emph{no termina}
        sobre $w$.
\end{itemize}

\begin{definicion}[Turing-reconocible]
\label{def:reconocible}
Un lenguaje $L\subseteq\Sstar$ es \emph{Turing-reconocible} (o,
equivalentemente, \emph{recursivamente enumerable}, o simplemente
\emph{reconocible}) si existe una MT $M$ tal que $\lang(M)=L$.
\end{definicion}

\begin{definicion}[Decidible]
\label{def:decidible}
Una MT es un \emph{decidor} si \emph{halta} en toda entrada (nunca cicla).
Un lenguaje $L\subseteq\Sstar$ es \emph{Turing-decidible} (o simplemente
\emph{decidible}, o \emph{recursivo}) si existe un decidor $M$ tal que
$\lang(M)=L$.
\end{definicion}

\begin{intuicion}
Reconocer un lenguaje es ``poder responder s\'i con certeza'' cuando la
respuesta es s\'i; pero, si la respuesta es no, podemos no obtener nunca esa
respuesta (la m\'aquina cicla). Decidir un lenguaje es algo m\'as fuerte:
``poder responder s\'i o no, siempre, en tiempo finito''.
\end{intuicion}

\subsection{Relaciones entre las clases}

\begin{proposicion}[Inclusi\'on]
Todo lenguaje decidible es reconocible: si $M$ es un decidor para $L$, en
particular $\lang(M)=L$.
\end{proposicion}

La inclusi\'on inversa es \emph{falsa}: existen lenguajes reconocibles pero
no decidibles. El ejemplo paradigm\'atico es $\ATM=\{\enc{M,w}\mid M
\text{ es MT y } M \text{ acepta }w\}$, que estudiaremos en el Pr\'actico 02.

\begin{ejemplo}[Reconocible no decidible: validez de primer orden]
Sea $L=\{\varphi : \,\vdash\varphi\}$, el conjunto de las formulas
v\'alidas (teoremas) de la l\'ogica de primer orden. Existe una MT que
reconoce $L$:
\begin{enumerate}
  \item En la cinta empieza una f\'ormula $\varphi$ codificada.
  \item La MT genera, una a una, todas las posibles demostraciones formales.
  \item Cuando aparece una demostraci\'on de $\varphi$, acepta.
\end{enumerate}
Si $\varphi$ es demostrable, eventualmente aparecer\'a; si no lo es, la
m\'aquina nunca termina. Se demuestra (cl\'asico, G\"odel/Turing) que $L$
\emph{no} es decidible: no existe un algoritmo que termine siempre con la
respuesta correcta sobre la validez.
\end{ejemplo}

\begin{observacion}
La diferencia entre reconocer y decidir es \emph{exactamente} la diferencia
entre ``s\'i o no, eventualmente'' y ``s\'i o no, en tiempo finito
garantizado''. Esa diferencia es responsable de fen\'omenos como el
problema de la detenci\'on (Halting Problem).
\end{observacion}

\subsection{Reconocibilidad por enumeradores}

Los \emph{enumeradores} dan una caracterizaci\'on alternativa muy \'util de
los lenguajes reconocibles. Un enumerador es una MT con dos cintas: una de
trabajo y otra (\emph{impresora}) en la que va \emph{imprimiendo} cadenas.

\begin{definicion}[Enumerador]
\label{def:enumerador}
Un \emph{enumerador} es una MT con dos cintas
$E=(Q,\Sigma,\Gamma_0,\Gamma_1,q_0,q_{\text{init}},\qaccept,\qreject,q_{\text{print}})$,
donde $\Gamma_0$ es el alfabeto de la cinta de trabajo, $\Gamma_1$ el de la
cinta de impresi\'on, y $q_{\text{print}}$ un estado distinguido tal que,
al alcanzarlo, $E$ ``imprime'' el contenido actual de la cinta de
impresi\'on (lo emite como salida) y la limpia para empezar a imprimir la
siguiente cadena.

El \emph{lenguaje enumerado} por $E$ es
\(L(E) = \{\,s : E \text{ imprime } s \text{ en alg\'un momento}\,\}.\)
\end{definicion}

\begin{teorema}[Enumeradores caracterizan reconocibilidad, Sipser 3.21]
\label{teo:enumerador}
Un lenguaje $L$ es Turing-reconocible si y solo si existe un enumerador $E$
con $L(E)=L$.
\end{teorema}

\begin{proof}[Idea de la prueba]
\textbf{($\Leftarrow$)} Si $E$ enumera $L$, construimos una MT $M$ que
reconoce $L$ as\'i: con entrada $w$, simula $E$ paso a paso; cada vez que
$E$ imprime una cadena $s$, $M$ compara $s$ con $w$; si son iguales, $M$
acepta. Si $w\in L$, eventualmente $E$ imprime $w$ y $M$ lo detecta.

\textbf{($\Rightarrow$)} Si $M$ reconoce $L$, se construye un enumerador
$E$ as\'i: enumera todas las cadenas $s_1,s_2,\dots$ de $\Sstar$ y, en la
iteraci\'on $i$, simula $M$ por $i$ pasos sobre cada $s_j$ con $j\le i$.
Cuando $M$ acepta alguna $s_j$, $E$ la imprime. Esto evita ciclar en una
$s_j$ que cicla en $M$: usa una estrategia de \emph{dovetailing}.
\end{proof}

\begin{intuicion}
Un enumerador puede imprimir las cadenas en cualquier orden y con
repeticiones. La caracterizaci\'on dice que ``poder reconocer'' (responder
s\'i a las cadenas del lenguaje) es lo mismo que ``poder enumerarlas''.
\end{intuicion}

\bigskip\hrule\bigskip

% =============================================================================
\section{Niveles de descripci\'on de una MT}
\label{sec:niveles}

Sipser (Secci\'on 3.3) distingue \textbf{tres niveles} para describir
m\'aquinas de Turing. La elecci\'on del nivel depende de qu\'e tan
detallado deba ser el an\'alisis.

\begin{enumerate}
  \item \textbf{Descripci\'on formal:} se especifican expl\'icitamente los
        siete componentes de la 7-tupla, en particular la funci\'on $\delta$
        como tabla o diagrama de estados. Es el nivel m\'as detallado y
        riguroso, pero impr\'actico para m\'aquinas grandes.
  \item \textbf{Descripci\'on a nivel implementaci\'on:} se da en lenguaje
        natural, indicando c\'omo la m\'aquina mueve el cabezal y modifica
        la cinta paso a paso (escribe, marca, recorre, vuelve, etc.), sin
        listar los estados ni la $\delta$. Es el nivel preferido para los
        ejercicios.
  \item \textbf{Descripci\'on de alto nivel:} se da el algoritmo a nivel
        l\'ogico, ignorando c\'omo se gestionan la cinta y el cabezal. Es
        el nivel del ``pseudoc\'odigo'' y se justifica por la
        \emph{tesis de Church-Turing}.
\end{enumerate}

\begin{intuicion}
Es como programar: el nivel formal es ensamblador, el de implementaci\'on es
C, y el de alto nivel es Python o pseudoc\'odigo. Todos describen el mismo
algoritmo, pero a distinto detalle.
\end{intuicion}

\bigskip\hrule\bigskip

% =============================================================================
\section{Ejemplos cl\'asicos}
\label{sec:ejemplos}

Repasamos algunas MT de referencia que aparecen en Sipser y en las l\'aminas
de la c\'atedra. Estos ejemplos ilustran t\'ecnicas que se aplican
directamente en los ejercicios.

\subsection{$M_1$ y el lenguaje $\{w\#w\}$ (Sipser 3.9)}

Sea $A=\{w\#w \mid w\in\{0,1\}^*\}$. La MT $M_1$ que decide $A$ funciona en
``zigzag'': para cada s\'imbolo a la izquierda del \texttt{\#}, marca el
correspondiente a su derecha. Si encuentra discrepancia, rechaza; si los
consume todos sin discrepancias, acepta.

T\'ecnica clave: \textbf{marcar celdas} con simbolos auxiliares
($x,y\in\Gamma\setminus\Sigma$) para recordar qu\'e ya se proces\'o.

\subsection{$M_2$ y el lenguaje $\{0^{2^n}\mid n\ge 0\}$ (Sipser 3.7)}

$M_2$ decide $A=\{0^{2^n}\}$, las cadenas de ceros cuya longitud es
potencia de $2$. La estrategia:
\begin{enumerate}
  \item recorrer la cinta tachando uno de cada dos $0$s (con $x$);
  \item si tras la pasada queda un \'unico $0$, aceptar;
  \item si la cantidad de $0$s sin tachar es impar y mayor que $1$, rechazar;
  \item si es par, volver al inicio y repetir.
\end{enumerate}
Cada iteraci\'on divide a la mitad la cantidad de $0$s; si la longitud no
es potencia de $2$, en alguna iteraci\'on aparecer\'a paridad impar.

T\'ecnica clave: \textbf{contador de paridad implementado en estados}
($q_3,q_4$ alternan).

\subsection{$M_3$ y el lenguaje $\{a^ib^jc^k : i\cdot j = k\}$ (Sipser 3.11)}

$M_3$ usa la estrategia: por cada $a$, ``shuttlea'' entre los $b$ y los $c$
tachando uno de cada para emparejarlos en cantidad $j$, y as\'i descontar
$j$ veces de los $c$ por cada $a$. Si al consumir todas las $a$ los $c$
quedaron exactamente vac\'ios, acepta. Esta es la t\'ecnica que extenderemos
en el Ejercicio 3 del pr\'actico.

\bigskip\hrule\bigskip

% =============================================================================
\section{Variantes del modelo y robustez}
\label{sec:variantes}

Una caracter\'istica notable de las MT es su \textbf{robustez}: muchas
modificaciones ``naturales'' del modelo no cambian la clase de lenguajes
reconocibles. Esto justifica la tesis de Church-Turing y nos permite
elegir, en cada situaci\'on, la variante m\'as conveniente.

\subsection{Cinta doblemente infinita}

\begin{definicion}[MT con cinta doblemente infinita]
Una MT con cinta doblemente infinita es id\'entica al modelo est\'andar,
salvo que su cinta es infinita en \emph{ambos} sentidos (las celdas est\'an
indexadas por $\mathbb Z$ en lugar de $\mathbb N$). Inicialmente la entrada
$w_0\cdots w_{n-1}$ ocupa las celdas $0,1,\dots,n-1$ y el resto contiene
blancos.
\end{definicion}

\begin{teorema}[Equivalencia, Problema 3.11 de Sipser]
\label{teo:doblemente-infinita}
Las MT con cinta doblemente infinita reconocen exactamente los mismos
lenguajes que las MT est\'andar.
\end{teorema}

\begin{proof}[Idea de la prueba]
La inclusi\'on est\'andar $\subseteq$ doblemente-infinita es trivial (basta
no usar la zona de la izquierda). La rec\'iproca usa la t\'ecnica de
``doblar la cinta'': se codifican las celdas $i\in\mathbb Z$ en celdas
$e(i)\in\mathbb N$ con
\[
  e(i) = \begin{cases} 2i & i\ge 0,\\ -2i-1 & i<0,\end{cases}
\]
de modo que las posiciones no negativas van a celdas pares y las negativas
a impares. Una marca distinguida indica la posici\'on simulada del cabezal.
\end{proof}

Este teorema corresponde directamente al Ejercicio~4 del pr\'actico.

\subsection{Stay put en lugar de Left}

Sipser (Problema 3.13) considera una variante donde el cabezal puede
moverse a la derecha o quedarse en su lugar (pero no ir a la izquierda).
\textbf{Esta variante no es equivalente} a la MT est\'andar: es estrictamente
m\'as d\'ebil. Es un buen contraste para entender que no toda variaci\'on
del modelo preserva el poder.

\subsection{MT con muchas cintas (multicinta)}

\begin{definicion}[Multicinta]
Una MT con $k$ cintas es una tupla
\[
  M = (Q, \Sigma, \{\Gamma_i\}_{0\le i<k}, \delta, q_0, \qaccept, \qreject),
\]
donde cada cinta $i$ tiene su propio alfabeto $\Gamma_i$ y su propio cabezal,
y la funci\'on de transici\'on es
\[
  \delta : Q \times \Gamma_0 \times \cdots \times \Gamma_{k-1}
        \to Q \times \Gamma_0 \times \cdots \times \Gamma_{k-1} \times \{L,R\}^k.
\]
La entrada se ubica en la cinta $1$; el resto comienzan en blanco.
\end{definicion}

\begin{teorema}[Multicinta $\equiv$ una cinta, Sipser 3.13]
\label{teo:multi-equiv}
Toda MT con $k$ cintas tiene una MT equivalente con una sola cinta.
\end{teorema}

\begin{proof}[Idea de la prueba]
Sea $M$ con $k$ cintas. Construimos $S$ con una sola cinta que codifica las
$k$ cintas en bloques separados por un nuevo s\'imbolo \texttt{\#}:
\[
  \texttt{\#}\,u_1\,\texttt{\#}\,u_2\,\texttt{\#}\,\cdots\,\texttt{\#}\,u_k\,\texttt{\#}.
\]
Para indicar las posiciones de los $k$ cabezales, $S$ marca un s\'imbolo de
cada bloque con un punto (\emph{dotted}). Para simular un paso de $M$:
\begin{enumerate}
  \item $S$ recorre toda la cinta y memoriza en su control finito los $k$
        s\'imbolos marcados.
  \item Calcula $\delta_M$ con esos s\'imbolos.
  \item Vuelve a recorrer la cinta y, en cada bloque, escribe el nuevo
        s\'imbolo y mueve la marca.
  \item Si una marca debe avanzar m\'as all\'a del fin de su bloque, $S$
        ``expande'' el bloque desplazando todo a la derecha.
\end{enumerate}
\end{proof}

\begin{corolario}
$L$ es Turing-reconocible si y s\'olo si lo reconoce alguna MT multicinta.
\end{corolario}

\subsection{No determinismo}

\begin{definicion}[MT no determinista, NTM]
Una NTM se obtiene cambiando la signatura de $\delta$:
\[
  \delta : Q\times\Gamma \to \Pow(Q\times\Gamma\times\{L,R\}).
\]
En cada configuraci\'on, hay un conjunto finito de transiciones posibles.
La m\'aquina \emph{acepta} $w$ si \emph{alguna} rama de su \'arbol de
ejecuci\'on alcanza $\qaccept$.
\end{definicion}

\begin{teorema}[Determinismo $\equiv$ no determinismo, Sipser 3.16]
Para toda NTM existe una MT determinista equivalente.
\end{teorema}

\begin{proof}[Idea de la prueba]
Se simula la NTM con un \textbf{BFS por capas} sobre el \'arbol de
configuraciones: la MT determinista visita las configuraciones de
profundidad $1$, luego de profundidad $2$, etc. Si alguna rama acepta, el
BFS la encontrar\'a en tiempo finito. Si ninguna acepta, el BFS contin\'ua
indefinidamente (esto refleja que la NTM puede tener ejecuciones que no
terminan).
\end{proof}

\begin{observacion}
El no determinismo no agrega poder de reconocimiento, pero a veces
\emph{simplifica} el dise\~no de una MT (especialmente en la teor\'ia de
complejidad).
\end{observacion}

\subsection{Restricciones que s\'i cambian el poder}

No todas las variaciones preservan poder. Por ejemplo, una MT que
\textbf{no puede escribir sobre el input} (Problema 3.20 de Sipser, y
Ejercicio~6 del pr\'actico) reconoce \textbf{solo lenguajes regulares}: el
modelo se reduce a un aut\'omata finito de dos v\'ias (2DFA), que es
equivalente a un DFA por el cl\'asico teorema de Rabin--Scott /
Shepherdson.

\bigskip\hrule\bigskip

% =============================================================================
\section{La tesis de Church-Turing}
\label{sec:church-turing}

\begin{teorema}[Tesis de Church-Turing]
La noci\'on intuitiva de \emph{algoritmo} coincide con la noci\'on formal
de \emph{computable por una MT}.
\end{teorema}

No es un teorema en sentido estricto (la noci\'on intuitiva no es
formalizable); es una \emph{tesis cient\'ifica} sostenida por evidencia
acumulativa: todas las formalizaciones razonables propuestas (MT,
$\lambda$-c\'alculo de Church, funciones recursivas de G\"odel-Kleene,
m\'aquinas de Post, RAM, etc.) resultan \textbf{equivalentes en poder}.

\begin{intuicion}
La tesis es la que justifica que, en niveles altos de descripci\'on, podamos
``hablar de algoritmos'' sin referirnos expl\'icitamente a una MT: si algo
se puede describir como un procedimiento finito y bien definido, entonces
existe una MT que lo computa. Esto es lo que permite que el Ejercicio~3 del
pr\'actico se resuelva con descripciones tipo Sipser sin escribir 100
l\'ineas de $\delta$.
\end{intuicion}

\bigskip\hrule\bigskip

% =============================================================================
\section{Cardinalidad y existencia de lenguajes no computables}
\label{sec:cardinalidad}

Esta secci\'on resume el contenido de \texttt{Diag.pdf} y motiva la
existencia de lenguajes no reconocibles --- y por ende, no decidibles --- por
argumentos de cardinalidad.

\subsection{Cardinalidad y conjuntos contables}

\begin{definicion}[Igualdad de cardinalidad]
Dos conjuntos $A$ y $B$ tienen la misma cardinalidad, $A\cong B$, si existe
una funci\'on biyectiva $f:A\to B$.
\end{definicion}

\begin{definicion}[Conjunto contable]
$A$ es \emph{contable} si $A\cong\mathbb N$. Equivalentemente, los elementos
de $A$ pueden listarse en una sucesi\'on $a_1,a_2,a_3,\dots$
\end{definicion}

\begin{ejemplo}
$\mathbb Z$ es contable, v\'ia
$f(n)=\lceil n/2\rceil$ si $n$ impar, $f(n)=-n/2$ si $n$ par.
\end{ejemplo}

\begin{ejemplo}
$\mathbb Q$ es contable: la enumeraci\'on diagonal de $(p,q)\in\mathbb N\times\mathbb N^+$ recorre todos los pares y, eliminando los no
reducidos, lista todos los racionales positivos. Combinando con los
negativos y el cero se obtiene una enumeraci\'on de $\mathbb Q$.
\end{ejemplo}

\subsection{$\mathbb R$ no es contable: el m\'etodo de la diagonal}

\begin{teorema}[Cantor]
$\mathbb R$ no es contable.
\end{teorema}

\begin{proof}
Basta mostrar que $[0,1)$ no es contable. Supongamos por absurdo que
$[0,1)\cong\mathbb N$ y enumeremos los elementos de $[0,1)$ por su expansi\'on
decimal:
\[
  r_1 = 0.d_1^1 d_2^1 d_3^1 \cdots
  \quad
  r_2 = 0.d_1^2 d_2^2 d_3^2 \cdots
  \quad
  r_3 = 0.d_1^3 d_2^3 d_3^3 \cdots
  \quad \cdots
\]
Construimos un n\'umero $r=0.\bar d^1 \bar d^2 \bar d^3 \cdots$ con
\[
  \bar d^i = \begin{cases} 1 & d_i^i \neq 1,\\ 0 & d_i^i = 1.\end{cases}
\]
Por construcci\'on $r$ difiere de $r_i$ en la $i$-\'esima cifra decimal,
para todo $i$. Por lo tanto $r\notin\{r_1,r_2,\dots\}$, contradiciendo que
la enumeraci\'on era completa.
\end{proof}

Este es el \textbf{m\'etodo de la diagonal de Cantor}, una herramienta clave
para probar resultados de no computabilidad.

\subsection{Existen lenguajes no computables}

\begin{teorema}
Existen lenguajes $L\subseteq\{0,1\}^*$ que no son reconocibles por
ninguna MT.
\end{teorema}

\begin{proof}
\begin{itemize}
  \item El conjunto de descripciones de MT $\{\enc{M}: M\text{ es MT}\}$
        es un subconjunto contable de $\{0,1\}^*$. Como $\{0,1\}^*$ es
        contable, su cardinalidad es $\aleph_0$.
  \item El conjunto de \emph{lenguajes} sobre $\{0,1\}^*$ es
        $\Pow(\{0,1\}^*)$. Por el teorema de Cantor, $\Pow(\{0,1\}^*)$
        tiene cardinalidad $2^{\aleph_0}>\aleph_0$.
\end{itemize}
Cada MT reconoce un lenguaje, as\'i que la cantidad de lenguajes
reconocibles es a lo sumo $\aleph_0$. Como hay $2^{\aleph_0}$ lenguajes en
total, debe haber lenguajes no reconocibles (de hecho, ``casi todos'' lo
son).
\end{proof}

\begin{intuicion}
El argumento es \emph{no constructivo}: dice que existen lenguajes no
reconocibles pero no exhibe ninguno. La construcci\'on de un ejemplo
expl\'icito (como $\overline{\ATM}$, Cap\'itulo 4 de Sipser) requiere
diagonalizaci\'on como en la prueba de Cantor.
\end{intuicion}

\bigskip\hrule\bigskip

% =============================================================================
\section{El problema de la detenci\'on (preview)}
\label{sec:halting}

Aunque pertenece propiamente al Pr\'actico 02, es \'util introducir el
problema de la detenci\'on para entender por qu\'e existe la distinci\'on
entre reconocible y decidible.

\begin{definicion}
\(\ATM = \{\enc{M,w} \mid M \text{ es una MT que acepta } w\}.\)
\end{definicion}

\begin{teorema}[Indecibilidad de $\ATM$, Turing 1936]
$\ATM$ no es decidible.
\end{teorema}

\begin{proof}[Idea (diagonalizaci\'on)]
Supongamos que existe $H$ que decide $\ATM$. Construimos una nueva MT $D$
que, con entrada $\enc{M}$:
\begin{enumerate}
  \item simula $H$ sobre $\enc{M,\enc{M}}$;
  \item si $H$ acepta, $D$ rechaza; si $H$ rechaza, $D$ acepta.
\end{enumerate}
Ahora corremos $D$ con entrada $\enc{D}$:
\begin{itemize}
  \item Si $D$ acepta $\enc{D}$, entonces $H$ rechaz\'o $\enc{D,\enc{D}}$,
        es decir $D$ no acepta $\enc{D}$. Contradicci\'on.
  \item Si $D$ rechaza $\enc{D}$, entonces $H$ acept\'o $\enc{D,\enc{D}}$,
        es decir $D$ acepta $\enc{D}$. Contradicci\'on.
\end{itemize}
Por lo tanto, $H$ no existe.
\end{proof}

\begin{observacion}
$\ATM$ \emph{s\'i} es reconocible: la m\'aquina universal $U$ que simula
$M$ con entrada $w$ y acepta si $M$ acepta es un reconocedor para $\ATM$.
Por lo tanto $\ATM$ es un lenguaje \textbf{reconocible pero no decidible},
ejemplo concreto de la diferencia entre las dos clases.
\end{observacion}

\bigskip\hrule\bigskip

% =============================================================================
\section{Subconjuntos infinitos decidibles}
\label{sec:subconjunto}

Este resultado conecta directamente con el Ejercicio~5 del pr\'actico.

\begin{teorema}[Sipser, Problema 3.19]
\label{teo:sub-decidible}
Todo lenguaje infinito Turing-reconocible tiene un subconjunto infinito
decidible.
\end{teorema}

\begin{proof}[Esquema]
Sea $A$ infinito y reconocible. Por el Teorema~\ref{teo:enumerador} existe
un enumerador $E$ para $A$. Fijamos el orden \emph{shortlex} sobre $\Sstar$
(longitud, luego lexicogr\'afico). Construimos inductivamente:
\begin{itemize}
  \item $b_1$: primera cadena que imprime $E$.
  \item $b_{i+1}$: primera cadena que imprime $E$ con $b_{i+1}>b_i$ en
        shortlex.
\end{itemize}
Como $A$ es infinito, siempre existe alguna cadena de $A$ mayor que $b_i$,
y como $E$ enumera $A$, eventualmente la imprime. La sucesi\'on
$b_1<b_2<\cdots$ es estrictamente creciente, infinita, y est\'a contenida
en $A$.

Sea $B=\{b_1,b_2,\dots\}$. Para decidir $B$, dada $x$, generamos $b_1,b_2,\dots$ hasta encontrar el primer $b_k\ge x$:
\begin{itemize}
  \item si $b_k=x$, aceptamos;
  \item si $b_k>x$, rechazamos.
\end{itemize}
Bajo shortlex, hay finitas cadenas $\le x$, as\'i que en finitos pasos
encontramos $b_k\ge x$. Luego $B$ es decidible.
\end{proof}

\bigskip\hrule\bigskip

% =============================================================================
\section{T\'ecnicas para resolver el Pr\'actico 01}
\label{sec:tecnicas}

Resumimos las herramientas conceptuales y t\'ecnicas que aparecen una y
otra vez en los ejercicios.

\subsection*{T1. Convencion para escribir configuraciones}

Para el Ejercicio~2 (configuraciones de $M_2$), basta aplicar paso a paso
la regla $C\to_M C'$ usando la tabla de $\delta$. Una configuraci\'on
$u\,q\,v$ codifica simult\'aneamente cinta, posici\'on del cabezal y estado.
El blanco $\blank$ s\'olo aparece a la derecha del segmento ``activo'' de la
cinta o como sentinela del extremo izquierdo.

\subsection*{T2. Marcado y emparejamiento}

Es la t\'ecnica central del Ejercicio~3 (deciders para
$\{w:\#_0(w)=\#_1(w)\}$ y $\{w:\#_0(w)=2\#_1(w)\}$). Se usan s\'imbolos
auxiliares (cl\'asicamente $x$ para ``$0$ usado'' e $y$ para ``$1$ usado'')
para emparejar simbolos del input. Cada iteraci\'on toma un simbolo no
marcado y busca su(s) contraparte(s); termina cuando ya no hay simbolos
sin marcar.

Plantilla:
\begin{itemize}
  \item Volver al inicio.
  \item Buscar el primer simbolo no marcado.
  \item Si no hay, aceptar (o pasar a una fase final de verificaci\'on).
  \item Si hay, marcarlo y buscar la contraparte. Si no existe, rechazar.
  \item Reiniciar.
\end{itemize}

\subsection*{T3. Simulaci\'on de variantes}

Para el Ejercicio~4 (cinta doblemente infinita) y para los algoritmos de
traducci\'on entre modelos:

\begin{itemize}
  \item \textbf{Doblemente infinita $\to$ est\'andar:} usar la codificaci\'on
        $e:\mathbb Z\to\mathbb N$, $e(i)=2i$ si $i\ge 0$, $e(i)=-2i-1$ si
        $i<0$.
  \item \textbf{Est\'andar $\to$ doblemente infinita:} reservar un marcador
        de extremo izquierdo $\vdash$ y nunca cruzarlo.
  \item \textbf{Multicinta $\to$ una cinta:} codificar las $k$ cintas en
        bloques separados por \texttt{\#}, con marcas para los $k$
        cabezales.
  \item \textbf{Una cinta $\to$ multicinta:} usar s\'olo la cinta 1 y dejar
        las restantes en blanco.
\end{itemize}

\subsection*{T4. Enumeraci\'on en orden creciente}

Para el Ejercicio~5: dado un enumerador $E$ de $A$, extraer la sucesi\'on
estrictamente creciente $b_1<b_2<\cdots$ tomando, sucesivamente, la primera
cadena impresa que sea mayor que la anterior. Decidir el subconjunto $B$
generado por esa sucesi\'on usando la finitud de $\{y\le x\}$ bajo
shortlex.

\subsection*{T5. Reducci\'on a un modelo m\'as d\'ebil}

Para el Ejercicio~6: si una MT no puede escribir sobre el input, su
comportamiento sobre la zona del input se reduce a movimientos del cabezal
sin alterar la cinta. Esto la convierte en un \textbf{2DFA}, y por
Rabin--Scott / Shepherdson los lenguajes reconocidos por 2DFA son los
regulares.

\bigskip\hrule\bigskip

% =============================================================================
\section{Resumen final}
\label{sec:final}

Hemos cubierto los conceptos centrales para abordar el Pr\'actico 01:

\begin{itemize}
  \item La \textbf{definici\'on formal} de MT como 7-tupla
        $(Q,\Sigma,\Gamma,\delta,q_0,\qaccept,\qreject)$ y la noci\'on de
        \emph{configuraci\'on} y \emph{c\'omputo}.
  \item La distinci\'on entre \textbf{lenguajes reconocibles} (existe MT
        que los acepta) y \textbf{decidibles} (existe decidor, MT que halta
        siempre). La inclusi\'on $\dec\subsetneq\rec$ es estricta.
  \item Las \textbf{descripciones} de MT en tres niveles (formal,
        implementaci\'on, alto nivel), justificadas por la tesis de
        Church-Turing.
  \item Las \textbf{variantes} principales del modelo: cinta doblemente
        infinita, multicinta, no determinista, enumeradores. Todas
        \emph{equivalentes} en poder; pero peque\~nas restricciones (como
        ``no escribir sobre el input'') s\'i pueden reducir el poder a la
        clase regular.
  \item La \textbf{tesis de Church-Turing}, que identifica algoritmo con
        c\'omputo por MT.
  \item La \textbf{existencia de lenguajes no computables} (Cantor) y un
        \textbf{ejemplo concreto} indecidible reconocible: $\ATM$.
  \item El \textbf{teorema sobre subconjuntos infinitos decidibles} y las
        \textbf{t\'ecnicas} para los ejercicios.
\end{itemize}

\paragraph{Hilo conductor.} La idea general detr\'as del Pr\'actico 01 es:
\emph{construir intuici\'on operativa con las MT, observar que peque\~nas
modificaciones del modelo no cambian su poder (robustez) salvo cuando
removemos alguna capacidad clave (como escribir), y entender por qu\'e existe
la distinci\'on fundamental entre reconocer y decidir, distinci\'on que
ser\'a el motor del Pr\'actico 02 sobre indecibilidad y reducciones.}

\bigskip
\hrule
\medskip
\begin{center}
\small
\emph{Fin del resumen te\'orico.} Para profundizar, leer el Cap\'itulo 3
completo de Sipser y revisar los ejercicios del Pr\'actico 01.
\end{center}

\end{document}
